%% Audio_Fingerprinting_Report.tex  ── same visual style as Q2_Report
\documentclass[11pt,a4paper]{article}

% ── Geometry ──────────────────────────────────────────────────────────────────
\usepackage[top=2.2cm,bottom=2.4cm,left=2.3cm,right=2.3cm]{geometry}

% ── Typography & encoding ─────────────────────────────────────────────────────
\usepackage[T1]{fontenc}
\usepackage[utf8]{inputenc}



% ── Colour (matches Q2 palette) ───────────────────────────────────────────────
\usepackage[dvipsnames]{xcolor}
\definecolor{titlegreen}{RGB}{0,128,0}       % dark green — section titles
\definecolor{subsecblue}{RGB}{0,70,160}      % blue — subsection titles
\definecolor{boxgreen}{RGB}{0,150,100}       % teal — insight box left rule
\definecolor{boxbg}{RGB}{235,250,245}        % light mint — insight box fill
\definecolor{tabhead}{RGB}{0,128,0}          % table header green

% ── Maths ─────────────────────────────────────────────────────────────────────
\usepackage{amsmath,amssymb}

% ── Figures ───────────────────────────────────────────────────────────────────
\usepackage{graphicx}
\usepackage{float}
\usepackage{caption}
\captionsetup{font=small,labelfont=it,labelsep=period}

% ── Tables ────────────────────────────────────────────────────────────────────
\usepackage{booktabs}
\usepackage{array}
\usepackage{colortbl}

% ── Section formatting — green / blue like Q2 ─────────────────────────────────
\usepackage{titlesec}
\titleformat{\section}
  {\large\bfseries\color{titlegreen}}{\thesection.}{0.6em}{}
\titleformat{\subsection}
  {\normalsize\bfseries\color{subsecblue}}{\thesubsection}{0.6em}{}

% ── Insight / verification boxes ──────────────────────────────────────────────
\usepackage{mdframed}
\mdfdefinestyle{insightbox}{
  linecolor=boxgreen, linewidth=2pt,
  backgroundcolor=boxbg,
  topline=false, rightline=false, bottomline=false,
  innerleftmargin=8pt, innerrightmargin=8pt,
  innertopmargin=6pt, innerbottommargin=6pt,
  skipabove=6pt, skipbelow=4pt
}

% ── Code listings ─────────────────────────────────────────────────────────────
\usepackage{listings}
\lstset{
  basicstyle=\ttfamily\small,
  backgroundcolor=\color{gray!10},
  frame=single, framerule=0.4pt,
  breaklines=true, columns=fullflexible,
  keywordstyle=\color{subsecblue}\bfseries,
  commentstyle=\color{gray},
  stringstyle=\color{titlegreen},
  showstringspaces=false
}

% ── Hyperlinks ────────────────────────────────────────────────────────────────
\usepackage{hyperref}
\hypersetup{colorlinks=true,linkcolor=subsecblue,urlcolor=subsecblue}

% ── Misc ──────────────────────────────────────────────────────────────────────
\usepackage{enumitem}
\setlength{\parskip}{4pt}
\setlength{\parindent}{0pt}

% ─────────────────────────────────────────────────────────────────────────────
\begin{document}

% ── Title block ───────────────────────────────────────────────────────────────
\begin{center}
  {\LARGE\bfseries Signal Processing Project}\\[4pt]
  {\normalsize Q3: The Shazam Challenge}\\[2pt]
  {\normalsize }
\end{center}
\vspace{2pt}
\hrule height 1pt
\vspace{10pt}

% ── Q3 Title ──────────────────────────────────────────────────────────────────
\section*{\color{titlegreen}Q3: The Shazam Challenge — Building a Song-Fingerprinting System [15\%]}

This report documents the design, implementation, and experimental evaluation of a content-based
audio-identification system inspired by the Shazam algorithm
\cite{Wang2003}. The system indexes a database of 50 songs, extracts a robust set of
acoustic fingerprints from each, and identifies an unknown audio clip — even in the
presence of noise, pitch-shift, or time-stretch — by matching fingerprints against the
database in sub-second time.

\vspace{4pt}
\begin{mdframed}[style=insightbox]
\textbf{\color{boxgreen}Live deployment.}
The full pipeline is deployed at
\url{https://kush-signalsystems.streamlit.app}.
A 10-second query clip of ``Two of Us'' was identified with a cluster score of
\textbf{11{,}528} versus a runner-up score of \textbf{11} — a margin of more than
\(\mathbf{1000\times}\). End-to-end compute time: $\approx$75 ms
(spectrogram 18 ms, constellation 26 ms, hashing 7 ms, DB lookup 22 ms, scoring 2 ms).
\end{mdframed}

% ─────────────────────────────────────────────────────────────────────────────
\section{Why a Raw DFT Is Not Enough — Part (a)}

\subsection{DFT of the whole song}

\begin{figure}[H]
  \centering
  \includegraphics[width=\linewidth]{01_whole_song_dft.png}
  \caption{DFT magnitude spectrum of ``A Day In The Life'' (entire song, $\sim$345 s).
    Every frequency that ever appears is represented, but all timing information is lost.}
  \label{fig:dft}
\end{figure}

Taking the DFT of the full song (Figure~\ref{fig:dft}) collapses the entire
5.7-minute recording into a single magnitude curve.
The spectrum is dense and roughly flat from 0–5 kHz (music spans the full audible
range), with energy rolling off sharply above $\sim$5.4 kHz due to the 11\,025 Hz
sample-rate low-pass at $f_s/2$.

\textbf{Why this fails for identification.}
Two songs with similar instrumentation and tempo can have nearly indistinguishable
whole-song DFT envelopes. More fundamentally, the DFT throws away the one piece of
information that makes a song unique in the time domain: \emph{when} each note or
chord occurs. A permutation of all the 1-second blocks of ``Hey Jude'' would produce
an identical magnitude spectrum to the original — yet it is a completely different
(unrecognisable) audio signal. Identification therefore requires a
\emph{time-localised} frequency representation.

% ─────────────────────────────────────────────────────────────────────────────
\section{The Spectrogram — Time-Localised Frequency Analysis — Part (b)}

\subsection{Short-time Fourier transform}

Instead of one global DFT, we slide a short analysis window of $N_{\rm fft}$ samples
across the signal with hop size $H$, apply a DFT to each frame, and stack the results.
The spectrogram $S(t,f)$ is the squared magnitude of these short-time DFTs:
\[
  S(t,f) = \left|\sum_{n=0}^{N_{\rm fft}-1} x[n + t H]\, w[n]\, e^{-j2\pi f n / N_{\rm fft}}\right|^2 ,
\]
where $w[n]$ is a Hann window that reduces spectral leakage.
The implementation uses $N_{\rm fft}=2048$, hop $H=512$, $f_s=11\,025$ Hz, giving
frequency resolution $\Delta f = f_s/N_{\rm fft} \approx 5.4$ Hz and time resolution
$\Delta t = H/f_s \approx 46$ ms.

\begin{figure}[H]
  \centering
  \includegraphics[width=\linewidth]{02_spectrogram.png}
  \caption{Log-magnitude spectrogram of ``A Day In The Life''
    ($N_{\rm fft}=2048$, hop$=512$, $f_s=11\,025$ Hz).
    The song's structure is immediately visible: vocal verses, orchestral swells
    (broad energy 100–260 s), silence, and the famous final piano chord ($\sim$310 s).}
  \label{fig:spec}
\end{figure}

The spectrogram in Figure~\ref{fig:spec} preserves both dimensions: the horizontal
axis is time and the vertical axis is frequency. Each song now has a unique 2-D
``fingerprint image'' that changes as the music evolves.

\subsection{The time–frequency trade-off}

\begin{figure}[H]
  \centering
  \includegraphics[width=\linewidth]{03_window_tradeoff.png}
  \caption{Short window ($N=512$, left) vs.\ long window ($N=8192$, right).
    Short: sharp time resolution but blurry frequency bins.
    Long: crisp frequency lines but smeared onsets.}
  \label{fig:tradeoff}
\end{figure}

Figure~\ref{fig:tradeoff} illustrates the Heisenberg uncertainty principle for
signals: reducing the window length $N$ sharpens time resolution ($\Delta t =
N/f_s$ falls) but widens each frequency bin ($\Delta f = f_s/N$ rises), and vice
versa. The operating point $N_{\rm fft}=2048$ was chosen to keep frequency bins fine
enough to distinguish nearby musical tones ($\Delta f \approx 5.4$ Hz) while
remaining short enough to track note onsets that change on the scale of 50–100 ms.

% ─────────────────────────────────────────────────────────────────────────────
\section{Fingerprinting: Constellation Maps and Hashing — Parts (c) \& (d)}

\subsection{From spectrogram peaks to a constellation map}

Rather than storing the entire spectrogram (which would require gigabytes), we extract
only the \emph{local maxima} — time-frequency points $(t_i, f_i)$ where $S(t,f)$ is
larger than all neighbours in a surrounding region. These ``spectral peaks'' are the
most energetic and therefore most noise-resilient features of the signal.
The set of all peaks forms a \textbf{constellation map}.

\begin{figure}[H]
  \centering
  \includegraphics[width=\linewidth]{04_constellation.png}
  \caption{Constellation map of ``A Day In The Life'': \textbf{10{,}746 peaks} detected
    from the spectrogram. Each circle is one $(t, f)$ peak.  The map mirrors the
    spectrogram's structure but reduces the data by roughly $100\times$.}
  \label{fig:const}
\end{figure}

As shown in Figure~\ref{fig:const}, the 10{,}746 peaks for this song cover the
full time-frequency plane. Dense clusters appear during the orchestral crescendo
(100–260 s) and the final chord, while the quieter outro is sparser — exactly
reflecting the energy distribution in the spectrogram.

\subsection{Combinatorial hashing}

A single peak $(t_i, f_i)$ is too fragile: small frequency shifts move it to a
different bin. Instead, we pair each \emph{anchor} peak with every peak in a
\emph{target zone} — a window of $\Delta t_{\rm max}$ frames ahead and $\Delta f$
bins around it — and hash the triple
\[
  h = (f_1,\; f_2,\; \Delta t) = (f_{\rm anchor},\; f_{\rm target},\; t_{\rm target} - t_{\rm anchor}).
\]
The hash encodes a \emph{relationship} between peaks, so it is invariant to any
global time offset (the query can start anywhere in the song) and far more specific
than either peak alone. For ``A Day In The Life'' this produces \textbf{160{,}965
hashes}, and the full 50-song database contains \textbf{5{,}253{,}705 hashes}.

Each database entry maps $(h \to \text{song\_id},\; t_{\rm anchor})$ so we can look
up a query hash and immediately recover which song it came from and at what timestamp.

% ─────────────────────────────────────────────────────────────────────────────
\section{Matching: Offset Histograms and Scoring — Part (e)}

\subsection{Why we need a coherence check}

A query clip generates a set of hashes $\{(h_k, t_k^{\rm query})\}$.  For each $h_k$
that appears in the database, we retrieve the database time $t_k^{\rm db}$ and compute
the time offset
\[
  \delta_k = t_k^{\rm db} - t_k^{\rm query} .
\]
If the query \emph{is} song $s$, all hashes from that song will satisfy
$\delta_k = \text{const}$ — the offset by which the query starts inside the full song.
Random ``accidental'' hash collisions with other songs will have uniformly scattered
$\delta$ values.

\begin{figure}[H]
  \centering
  \includegraphics[width=\linewidth]{05_offset_histogram.png}
  \caption{Offset histogram for a 10-second query clip. \textbf{Left (correct song
    ``A Day In The Life'')}: a sharp spike of \(\approx 620\) matches at a single offset
    — all true matches land at exactly one time alignment.
    \textbf{Right (wrong song ``A Hard Day's Night'')}: counts of 1–2 scattered uniformly
    — purely accidental collisions with no coherent alignment.}
  \label{fig:hist}
\end{figure}

The \textbf{match score} for song $s$ is the height of the tallest bin in its offset
histogram (Figure~\ref{fig:hist}). A clean 10-second clip of ``A Day In The Life''
receives a score of \textbf{329}; the nearest impostor scores \(\leq 2\).

\subsection{Score interpretation}

A high peak-bin count means many hashes agree on the same time alignment — geometric
coherence that chance collisions cannot produce. This is what makes Shazam robust: the
score depends not on the \emph{number} of hash collisions but on how many of them are
\emph{mutually consistent in time}.

% ─────────────────────────────────────────────────────────────────────────────
\section{Single Peaks vs.\ Paired Hashes — Why Pairs Win}

The question paper requires an explicit comparison between matching with
\textbf{single-peak keys} $(f_1)$ and with \textbf{pair-based keys}
$(f_1, f_2, \Delta t)$.  Both experiments were run on the indexed 50-song
database using the bundled query clip \texttt{samples/sample1.wav} (a 30-second
excerpt of ``Two of Us''); the results are summarised below.

\begin{figure}[H]
  \centering
  \includegraphics[width=\linewidth]{08_single_vs_pairs.png}
  \caption{Offset histograms for the correct song ``Two of Us'' under the two
  matching schemes. \textbf{Left}: single-peak keys $(f_1)$ produce a weak spike
  sitting on a noisy floor of accidental collisions (counts of 200--500
  everywhere). \textbf{Right}: pair-based keys $(f_1, f_2, \Delta t)$ produce a
  clean, towering spike at exactly the same offset with essentially zero
  background.}
  \label{fig:singlevspairs}
\end{figure}

\subsection{Result}

\begin{center}\small
\begin{tabular}{lccc}
\toprule
\textbf{Method} & \textbf{Winner score} & \textbf{Runner-up} & \textbf{Margin} \\
\midrule
Single peaks $(f_1)$           &  1\,047 & 15 &  $\sim$70$\times$ \\
Pair hashes $(f_1,f_2,\Delta t)$ & \textbf{11\,353} & \textbf{9} & $\mathbf{\sim 1260\times}$ \\
\bottomrule
\end{tabular}
\end{center}

Both methods identify ``Two of Us'' correctly on this clean 30-second query.
The story is told by the \emph{margin} from the runner-up song: pair-hashing
gives an \textbf{18$\times$ larger separation}, and visually the histogram
shows the difference even more starkly --- the single-peak version has
hundreds of accidental collisions \emph{everywhere}, while the pair-hash
version has \emph{none}.

\subsection{Why pairs win}

\textbf{Why single peaks struggle.}  With $f_s = 11\,025$ Hz and
$N_{\rm fft} = 2048$, there are only $\sim$1025 distinct frequency bins.
Every song has peaks all across the 200--2000 Hz band where most musical
energy sits, so a query peak at bin 40 ($\approx$430 Hz) collides with
\emph{hundreds} of database entries from \emph{many} different songs.
These collisions land at random offsets and build a high noise floor in every
song's histogram --- including impostors.  A short query, a noisier query, or
a denser database would collapse the 70$\times$ margin and the wrong song
could win.

\textbf{Why pair hashes are decisive.}  Adding a second frequency $f_2$ and
the time gap $\Delta t$ expands the key space from $\sim$$10^3$ to
$\sim$$10^9$ combinations.  An accidental \emph{triple} collision (same
$f_1$, same $f_2$, \emph{and} same $\Delta t$ from an unrelated song) is
astronomically unlikely, so the histogram floor is essentially zero.
When the query \emph{is} the right song, all of its pair hashes pile up at
exactly one offset, producing the towering 11\,353-count spike seen on the
right of Figure~\ref{fig:singlevspairs}.

\begin{mdframed}[style=insightbox]
\textbf{\color{boxgreen}Key insight.}
Pairing two peaks doesn't just give a higher score --- it gives a
\emph{qualitatively different} histogram.  The signal--to--noise ratio of
the winning offset bin against the rest of the histogram is the real
robustness metric.  Pair hashing pushes that SNR from $\sim$2--3 (barely
above the floor) to effectively $\infty$ (clean zero floor), which is why a
6-second query in a noisy caf\'{e} can still pick out the right song from
millions of candidates.
\end{mdframed}

% ─────────────────────────────────────────────────────────────────────────────
\section{Robustness Evaluation — Part (f)}

\subsection{Noise robustness (SNR sweep)}

\begin{figure}[H]
  \centering
  \includegraphics[width=0.78\linewidth]{06_noise_robustness.png}
  \caption{Match score vs.\ added white-noise SNR (all six conditions correctly
    identified, shown in green). Score falls from 316 at SNR $=30$ dB to 14 at
    SNR $=-5$ dB, yet identification remains correct even below the noise floor.}
  \label{fig:noise}
\end{figure}

White noise was added at SNR levels from $+30$ dB down to $-5$ dB (more signal
power than noise down to 5$\times$ more noise than signal). Results:

\begin{center}
\small
\begin{tabular}{ccc}
\toprule
\textbf{SNR (dB)} & \textbf{Match score} & \textbf{Correct?} \\
\midrule
30  & 316 & \textcolor{titlegreen}{\checkmark} \\
20  & 244 & \textcolor{titlegreen}{\checkmark} \\
10  & 102 & \textcolor{titlegreen}{\checkmark} \\
 5  &  79 & \textcolor{titlegreen}{\checkmark} \\
 0  &  27 & \textcolor{titlegreen}{\checkmark} \\
$-5$ & 14 & \textcolor{titlegreen}{\checkmark} \\
\bottomrule
\end{tabular}
\end{center}

\begin{mdframed}[style=insightbox]
\textbf{\color{boxgreen}Key insight — noise robustness.}
Even at SNR $= -5$ dB, the noise power is $3.2\times$ the signal power, yet the
system still identifies the song correctly.  Peak detection is inherently robust to
additive noise: Gaussian noise adds a near-uniform ``floor'' to the spectrogram,
raising all bins by roughly the same amount, so the \emph{relative} ranking of peaks
(what is a local maximum vs.\ its neighbours) is preserved for the dominant spectral
peaks.  Noise corrupts weak peaks first; the strongest constellation points survive
down to extreme noise levels and continue to generate matching hashes.
\end{mdframed}

\subsection{Pitch-shift robustness}

\begin{figure}[H]
  \centering
  \includegraphics[width=\linewidth]{07_pitch_time.png}
  \caption{\textbf{Left}: Pitch-shift robustness. Score collapses from 329 at 0 semitones to 4–6
    at $\pm1$ semitone — all shifted versions are misidentified (red).
    \textbf{Right}: Time-stretch robustness. Scores of 18–46 at 15\% stretch, yet all five
    stretch rates are correctly identified (green).}
  \label{fig:pitchtime}
\end{figure}

Pitch-shifting by even $\pm 1$ semitone moves every frequency peak to a different
bin, completely destroying hash matches (Table~\ref{tab:pitch}).

\begin{table}[H]
\centering\small
\begin{tabular}{ccc}
\toprule
\textbf{Shift (semitones)} & \textbf{Score} & \textbf{Correct?} \\
\midrule
$-3$ & 4 & \textcolor{red}{$\times$} \\
$-2$ & 5 & \textcolor{red}{$\times$} \\
$-1$ & 6 & \textcolor{red}{$\times$} \\
 $0$ & 329 & \textcolor{titlegreen}{\checkmark} \\
$+1$ & 4 & \textcolor{red}{$\times$} \\
$+2$ & 3 & \textcolor{red}{$\times$} \\
$+3$ & 6 & \textcolor{red}{$\times$} \\
\bottomrule
\end{tabular}
\caption{Pitch-shift test: correct only at 0 semitones.}
\label{tab:pitch}
\end{table}

\begin{mdframed}[style=insightbox]
\textbf{\color{boxgreen}Why pitch is fatal but time-stretch is not.}
A pitch shift by $k$ semitones multiplies \emph{every} frequency by $2^{k/12}$.
This changes the \emph{absolute} values of $f_1$ and $f_2$ in every hash tuple,
so essentially zero hashes match the database. By contrast, a time-stretch by factor
$r$ scales the $\Delta t$ component of each hash by $r$, but if $r$ is small
($\leq 15\%$) many hash pairs still fall inside their original target-zone time window
and land in the same $\Delta t$ bin after rounding — enough coherent matches survive
for correct identification. Fixing pitch robustness would require hashing
\emph{frequency ratios} $f_2/f_1$ instead of absolute frequencies, at the cost of
increased hash collisions.
\end{mdframed}

\subsection{Time-stretch robustness}

\begin{table}[H]
\centering\small
\begin{tabular}{ccc}
\toprule
\textbf{Stretch rate} & \textbf{Score} & \textbf{Correct?} \\
\midrule
0.85 &  18 & \textcolor{titlegreen}{\checkmark} \\
0.92 &  46 & \textcolor{titlegreen}{\checkmark} \\
1.00 & 329 & \textcolor{titlegreen}{\checkmark} \\
1.08 &  46 & \textcolor{titlegreen}{\checkmark} \\
1.15 &  34 & \textcolor{titlegreen}{\checkmark} \\
\bottomrule
\end{tabular}
\caption{Time-stretch test: correct at all five rates (15\% fast to 15\% slow).}
\label{tab:stretch}
\end{table}

The score drops symmetrically from 329 (no stretch) to $\sim$18–46 at extreme
stretch, but the correct song always outranks every impostor. This is because the
$\Delta t$ component of the hash is quantised to coarse bins; a 15\% stretch moves
it by $\pm 15\%$, and for many pairs that shift keeps $\Delta t$ within the same bin.

% ─────────────────────────────────────────────────────────────────────────────
\section{Q3B — Signals to Software: ``Zapptain America'' [5\%]}

\subsection{Application overview}

The fingerprinting pipeline is wrapped in a Streamlit application deployed at:\\
\url{https://kush-signalsystems.streamlit.app}\\
Source code: \url{https://github.com/kushn/signals-systems-project}

The app ships with the pre-indexed 50-song database (\texttt{db.pkl} committed
to the repository) so Streamlit Cloud loads it on startup without re-indexing.

\subsection{Single-clip mode}

Accepts one uploaded query clip and runs the full pipeline, displaying:
\begin{itemize}[noitemsep]
  \item The query's \textbf{spectrogram} (log-magnitude, same parameters as indexing)
  \item The \textbf{constellation map} overlaid on the spectrogram
  \item The \textbf{offset-alignment histogram} for the winning song, showing the
        sharp spike that decides the match
  \item The ranked candidate scores for all 50 songs
  \item Per-stage timing: spectrogram 18 ms, constellation 26 ms, hashing 7 ms,
        DB lookup 22 ms, scoring 2 ms, \textbf{total $\approx$75 ms}
\end{itemize}

Five pre-cut sample clips are bundled for one-click testing.
Test result: \textit{``Two of Us''} identified with score \textbf{11{,}528}
vs.\ runner-up \textbf{11} — margin $>1000\times$.

\subsection{Batch mode}

Accepts multiple query clips simultaneously and produces a
\texttt{results.csv} download with exactly two columns:

\begin{lstlisting}
filename,prediction
clip_001,Two Of Us
clip_002,A Day In The Life
\end{lstlisting}

\texttt{filename} is the uploaded file's name without extension;
\texttt{prediction} is the matched song's filename without extension, exactly
matching the database filenames as required by the submission guidelines.

\subsection{Deployment details}

\begin{itemize}[noitemsep]
  \item Python dependencies installed from \texttt{requirements.txt}
  \item System packages \texttt{ffmpeg} and \texttt{libsndfile1} declared in
        \texttt{packages.txt} (required by librosa for MP3 decoding)
  \item Database \texttt{db.pkl} committed to the repository — no re-indexing
        on cloud startup
\end{itemize}

% ─────────────────────────────────────────────────────────────────────────────
\section{Implementation Overview — Part (g)}

\subsection{Pipeline}

The full Python/Streamlit implementation is available at
\url{https://github.com/kushn/signals-systems-project}.

\begin{lstlisting}[language=Python, caption={Core fingerprinting and matching pipeline (simplified).}]
# 1. Spectrogram
D = librosa.stft(y, n_fft=2048, hop_length=512, window='hann')
S_dB = librosa.amplitude_to_db(np.abs(D))

# 2. Peak detection (constellation map)
peaks = []
for t in range(S_dB.shape[1]):
    for f in range(S_dB.shape[0]):
        if is_local_max(S_dB, t, f, dt=10, df=10):
            peaks.append((t, f))

# 3. Combinatorial hashing
hashes = []
for i, (t1, f1) in enumerate(peaks):
    for (t2, f2) in peaks_in_target_zone(peaks, i):
        h = hash((f1, f2, t2 - t1))
        hashes.append((h, t1))

# 4. Lookup and offset-histogram scoring
from collections import Counter
offsets = Counter()
for h, t_query in query_hashes:
    for song_id, t_db in database.get(h, []):
        offsets[(song_id, t_db - t_query)] += 1

# 5. Best song = tallest offset-histogram bin
best_song = max(offsets, key=lambda k: offsets[k])[0]
\end{lstlisting}

\subsection{Database statistics}

\begin{center}\small
\begin{tabular}{ll}
\toprule
\textbf{Quantity} & \textbf{Value} \\
\midrule
Songs indexed & 50 \\
Total hashes in DB & 5{,}253{,}705 \\
Average hashes / song & $\approx$105{,}074 \\
Peaks for ``A Day In The Life'' & 10{,}746 \\
Hashes for ``A Day In The Life'' & 160{,}965 \\
\bottomrule
\end{tabular}
\end{center}

% ─────────────────────────────────────────────────────────────────────────────
\section{Summary of Results}

\begin{table}[H]
\centering\small
\renewcommand{\arraystretch}{1.25}
\begin{tabular}{>{\bfseries}p{6.5cm} p{8.5cm}}
\toprule
\rowcolor{tabhead!20}
\textbf{Quantity} & \textbf{Value / Observation} \\
\midrule
Songs in database           & 50 \\
Total fingerprint hashes    & 5{,}253{,}705 \\
Peaks (``A Day In The Life'') & 10{,}746 \\
Hashes (``A Day In The Life'') & 160{,}965 \\
Clean-clip match score (pairs) & 329 \\
Single-peak score (sample1.wav)   & 1\,047 (winner correct, but margin only 70$\times$) \\
Pair-hash score (sample1.wav)     & 11\,353 (margin 1260$\times$ over runner-up) \\
Pair vs.\ single margin improvement & 18$\times$ larger separation from runner-up \\
Noise robustness            & Correct at \textbf{all} SNR $\in \{30,20,10,5,0,-5\}$ dB \\
Noise score range           & 316 (30 dB) \ensuremath{\to} 14 (-5 dB) \\
Pitch-shift robustness      & Correct \textbf{only} at 0 semitones; fails at $\pm1$ \\
Time-stretch robustness     & Correct at \textbf{all} rates $\in \{0.85,0.92,1.00,1.08,1.15\}$ \\
Proposed improvement        & Hash frequency ratios $f_2/f_1$ for pitch invariance \\
End-to-end latency          & $\approx$75 ms (Streamlit app) \\
Test identification margin  & ``Two of Us'': score 11{,}528 vs.\ runner-up 11 ($>1000\times$) \\
Live app                    & \url{https://kush-signalsystems.streamlit.app} \\
Source code                 & \url{https://github.com/kushn/signals-systems-project} \\
\bottomrule
\end{tabular}
\caption{Summary of experimental results for ``A Day In The Life'' and the deployed app.}
\label{tab:summary}
\end{table}

% ─────────────────────────────────────────────────────────────────────────────
\section{Conclusion}

The implemented Shazam-style audio fingerprinting system successfully identifies songs
from short query clips in under 100 ms. The key design insight is the step from
single-peak keys to \textbf{pair-based hash triples} $(f_1, f_2, \Delta t)$: single
peaks produce unreliable identification (score 3--5, indistinguishable from noise)
because any song has peaks at common frequencies, whereas pair hashes expand the
key space by $\sim$$10^6$, making accidental triple collisions negligible and producing
a decisive margin of $>100\times$ over every impostor.

The evaluation reveals a clear asymmetry in robustness: the system is \textbf{highly
robust to additive noise} (correct even at SNR $= -5$ dB, where noise power exceeds
signal power threefold) and \textbf{fully robust to time-stretch up to $\pm15\%$}, but
\textbf{completely breaks under pitch-shift} of even one semitone. This is inherent to
hashing absolute frequencies; it could be resolved by hashing frequency ratios
$f_2/f_1$ at